\documentclass[11pt,a4paper]{article}
\usepackage[utf8]{inputenc}
\usepackage[T1]{fontenc}
\usepackage[brazilian]{babel}
\usepackage{geometry}
\usepackage{enumitem}
\usepackage{booktabs}
\usepackage{amsmath}
\usepackage{graphicx}
\graphicspath{{../figures_global/}}
\usepackage{fancyhdr}
\usepackage[backend=biber, style=numeric, sorting=none]{biblatex}
\usepackage{hyperref}
\addbibresource{referencias.bib}

\geometry{margin=2.5cm, headheight=32pt}
\setlength{\parindent}{0pt}
\setlength{\parskip}{0.5em}

\pagestyle{fancy}
\fancyhf{}
\fancyhead[L]{\small\textbf{Faculdade de Ciências e Tecnologias em Engenharia}}
\fancyhead[R]{\includegraphics[height=0.85cm]{../figures_global/unb.png}}
\renewcommand{\headrulewidth}{0.4pt}
\fancyfoot[C]{\small Laboratório de Comunicações --- Ementa}
\fancyfoot[R]{\small \thepage}
\renewcommand{\footrulewidth}{0.4pt}
\fancypagestyle{plain}{%
  \fancyhf{}%
  \fancyhead[L]{\small\textbf{Faculdade de Ciências e Tecnologias em Engenharia}}%
  \fancyhead[R]{\includegraphics[height=0.85cm]{../figures_global/unb.png}}%
  \renewcommand{\headrulewidth}{0.4pt}%
  \fancyfoot[C]{\small Laboratório de Comunicações --- Ementa}%
  \fancyfoot[R]{\small \thepage}%
  \renewcommand{\footrulewidth}{0.4pt}%
}

\title{\textbf{Ementa}\\[0.3em]\large Laboratório de Comunicações}
\author{Faculdade de Ciências e Tecnologias em Engenharia --- FCTE/UnB}
\date{}

\begin{document}

\thispagestyle{plain}
\maketitle
\vspace{-1.0em}

\section{Identificação}

\begin{center}
\small
\begin{tabular}{@{}p{0.32\textwidth} p{0.60\textwidth}@{}}
\toprule
\textbf{Componente curricular} & Laboratório de Comunicações \\
\textbf{Natureza} & Obrigatória \\
\textbf{Créditos} & 2 créditos --- 1 teórico e 1 de extensão \\
\textbf{Carga horária total} & 30 horas \\
\textbf{Carga horária teórica} & 15 horas \\
\textbf{Carga horária de extensão} & 15 horas \\
\textbf{Modalidade} & Presencial, em laboratório \\
\textbf{Pré-requisito ou correquisito} & Princípios de Comunicação para Engenharia \\
\bottomrule
\end{tabular}
\end{center}

A carga horária de extensão corresponde a 50\% do componente e atende à curricularização da extensão prevista na Resolução CNE/CES 07/2018 e nas normas da UnB.

\section{Ementa}

Prática experimental de sistemas de comunicação em duas frentes complementares. Na primeira, implementação e análise de cadeias de modulação, recepção e sincronização utilizando rádio definido por software e o ambiente GNU Radio. Na segunda, medição de radiopropagação e caracterização de circuitos analógicos de radiofrequência, com avaliação de seu impacto sobre o desempenho dos sistemas de comunicação. As práticas articulam-se a um projeto de extensão de apoio técnico a equipes de competição da universidade.

\section{Objetivos}

\subsection*{Objetivo geral}

Desenvolver competências experimentais para implementar, medir e diagnosticar sistemas de comunicação, integrando o domínio de ferramentas de rádio definido por software com a caracterização de circuitos e enlaces reais, em benefício de demandas concretas de equipes e parceiros externos.

\subsection*{Objetivos específicos}

\begin{itemize}[noitemsep]
  \item Construir e analisar cadeias de modulação e demodulação analógica e digital em GNU Radio.
  \item Implementar blocos de recepção, incluindo filtragem casada e decisão.
  \item Aplicar técnicas de sincronização de portadora, temporização e quadro.
  \item Medir radiopropagação e verificar modelos de perda de percurso e orçamento de enlace.
  \item Caracterizar circuitos analógicos de radiofrequência quanto a ganho, ruído e linearidade.
  \item Relacionar as imperfeições de circuito ao desempenho medido do sistema de comunicação.
  \item Prestar apoio técnico a equipes de competição, transferindo os resultados de laboratório para uma aplicação real.
\end{itemize}

\section{Conteúdo Programático}

\begin{center}
\small
\begin{tabular}{|p{0.44\textwidth}|p{0.44\textwidth}|}
\hline
\textbf{Frente A: Modulação, Recepção e Sincronização em GNU Radio} & \textbf{Frente B: Radiopropagação e Caracterização de Circuitos Analógicos} \\
\hline
\parbox[t]{0.42\textwidth}{%
  \vspace{0.2em}
  \begin{itemize}[noitemsep, leftmargin=*, nosep, font=\small]
    \item Introdução ao rádio definido por software e ao GNU Radio
    \item Modulação e demodulação analógica AM e FM
    \item Modulação digital em banda-base e passa-banda
    \item Filtro casado e decisão na recepção
    \item Sincronização de portadora
    \item Sincronização de temporização de símbolo
    \item Detecção de preâmbulo e sincronização de quadro
    \item Recepção de sinais reais do ambiente
  \end{itemize}
  \vspace{0.2em}
} &
\parbox[t]{0.42\textwidth}{%
  \vspace{0.2em}
  \begin{itemize}[noitemsep, leftmargin=*, nosep, font=\small]
    \item Linhas de transmissão e casamento de impedância
    \item Medição de perda de percurso e modelos de propagação
    \item Diagrama de radiação de antenas
    \item Caracterização de amplificador de baixo ruído
    \item Não linearidades, compressão e intermodulação
    \item Filtros de radiofrequência e resposta em frequência
    \item Ruído de fase do oscilador e seu efeito na constelação
    \item Orçamento de enlace e verificação experimental
  \end{itemize}
  \vspace{0.2em}
} \\
\hline
\end{tabular}
\end{center}

A lista detalhada e o roteiro das práticas encontram-se no documento complementar de práticas do laboratório.

\section{Componente de Extensão}

A componente de extensão corresponde a 1 crédito e concretiza-se no apoio técnico a equipes de competição da universidade. As turmas colocam a instrumentação e os métodos do laboratório a serviço de uma necessidade real de comunicação da equipe, como um enlace de telemetria, a caracterização de uma antena ou o diagnóstico de um receptor.

\subsection*{Formas de interação com a comunidade externa}

\begin{itemize}[noitemsep]
  \item Diagnóstico e melhoria de enlaces de telemetria e telecomando de equipes de competição.
  \item Caracterização de antenas, cabos e módulos de radiofrequência usados pelas equipes.
  \item Apoio a demonstrações públicas, mostras e eventos de divulgação científica.
\end{itemize}

\subsection*{Produtos e registro da extensão}

Cada grupo entrega um laudo técnico das medições realizadas, as recomendações à equipe atendida e um registro da atividade conforme as normas de extensão da UnB, incluindo o cadastro da ação e o registro da carga horária dos estudantes.

\section{Metodologia}

O laboratório funciona por rotação em bancadas, com roteiros previamente disponibilizados e execução em grupos. A frente de GNU Radio é conduzida em rádio definido por software e a frente de radiofrequência utiliza instrumentação de bancada. A componente de extensão é planejada junto às equipes de competição, com definição de escopo, execução das medições e devolutiva dos resultados.

\section{Avaliação}

\subsection*{Componente prática}

A nota das práticas $N_P$ é a média dos relatórios das atividades experimentais das duas frentes:
\[
N_P = \frac{\text{Rel}_1 + \text{Rel}_2 + \cdots + \text{Rel}_n}{n}
\]

\subsection*{Componente de extensão}

A nota de extensão $N_E$ avalia o apoio prestado à equipe de competição por meio do laudo técnico, da devolutiva e do registro da atividade:
\[
N_E = 0{,}5\,\text{Laudo} + 0{,}3\,\text{Devolutiva} + 0{,}2\,\text{Registro}
\]

\subsection*{Nota final e aprovação}

\[
N_F = 0{,}5\,N_P + 0{,}5\,N_E
\]

\begin{itemize}[noitemsep]
  \item Exige-se aprovação simultânea nas práticas e na extensão, com $N_P \geq 5$ e $N_E \geq 5$.
  \item A aprovação final é obtida com $N_F \geq 5$.
  \item É obrigatória a frequência mínima de 75\% nas atividades de laboratório e de extensão.
\end{itemize}

\section{Referências}

\nocite{lathi2019modern, haykin2009communication, couch2013digital, proakis2001digital, carlson2002communication}
\printbibliography[heading=none]

\end{document}
